\section{The additional lead: local covering germs}
\label{sec:locallead}

We give ``sub-germ'' a precise geometric meaning: a connected finite
cover of the punctured strict henselian neighbourhood is a connected
component of the pullback of a finite global cover. A weaker meaning,
that the local cover is dominated by such a component, gives the same
criterion. The covers throughout this section are finite \'etale; a
ramified germ crossing the divisor is understood by its restriction to
the complement.

\begin{claim}{G.1}{Proved}
Let $i:H\to G$ be a continuous homomorphism of profinite groups.
The following assertions are equivalent:
\begin{enumerate}[label=(\roman*)]
\item $i$ is injective.
\item Every transitive finite continuous $H$-set is isomorphic to an
$H$-orbit in the restriction of a finite continuous $G$-set.
\item Every transitive finite continuous $H$-set is a quotient of an
$H$-orbit in the restriction of a finite continuous $G$-set.
\end{enumerate}
Thus the germ criterion in the prompt is true for the geometric local
tame group $H=\widehat{\Z}(1)^{I(x)}$ in characteristic zero, with
$G=\pi_1(U_{\overline K})$, provided that the local point is geometric
and the normal crossings divisor has its branches separated by strict
henselization. It is a criterion, not an assertion that local inertia
is always injective.
\end{claim}
\begin{proof}
Assume (i). A continuous injection from a compact space into a Hausdorff
space is a homeomorphism onto its closed image. We identify $H$ with
that image. A transitive finite continuous $H$-set has the form $H/V$
for an open subgroup $V$ of $H$. The subspace topology supplies an
open normal subgroup $N$ of $G$ with $N\cap H\subseteq V$.
Then $VN$ is a subgroup of $G$, because $N$ is normal, and it is open,
because it contains $N$. Moreover
\[
                    H\cap VN=V.
\]
Indeed $h=vn$ with $h,v\in H$ implies $n\in H\cap N\subseteq V$.
Consequently the orbit of $VN$ in $G/VN$ is exactly $H/V$.
This proves (ii), including its stronger exact-component conclusion.
Assertion (ii) implies (iii). If (iii) holds, every element of
$\ker i$ acts trivially on every finite transitive $H$-set. In
particular it acts trivially on $H/V$ for each open normal subgroup
$V$ of $H$, so belongs to $\bigcap_V V=1$. This proves (i).

For the geometric application choose a geometric generic base point
of the punctured strict henselian neighbourhood mapping to a geometric
base point of $U$. Pullback of finite \'etale covers is the restriction
functor for the induced continuous map $i$. The equivalence between
finite covers and finite continuous fundamental-group sets identifies
connected covers with transitive sets and connected components with
orbits. In a smooth normal crossings chart in characteristic zero,
Kummer covers of all orders in the branch parameters are cofinal:
Abhyankar's lemma removes the ramification, and a finite unramified
cover of a strictly henselian local ring is trivial. The local group
is therefore the stated product of Tate modules. Changing paths
conjugates $i$ in $G$, leaving every assertion invariant.
The precise local inputs are SGA~1 XIII, Propositions~5.1--5.2;
the abstract finite-cover criterion is also SGA~1 V, Corollary~6.8
\citep{E_SGA1}. These are statements about all finite covers and
the whole profinite group, not merely a chosen collection of cyclic
quotients.
\end{proof}

\begin{claim}{G.2}{Proved; explicit boundary cases}
The criterion holds for $U=\G_m^r\times\mathbb A^{d-r}$ at the
intersection of the coordinate boundary divisors in $\mathbb A^d$.
It fails for $\mathbb A^1\subset\mathbb P^1$ at infinity. At
$[0:0:1]\in\mathbb P^2$ with divisor $xy=0$, the local-to-global map
is
\[
 \widehat{\Z}(1)^2\longrightarrow\widehat{\Z}(1),
                  (a,b)\longmapsto a-b,
\]
and has a diagonal kernel.
\end{claim}
\begin{proof}
Over an algebraically closed characteristic-zero field the Kummer
covers of the coordinate torus give the first isomorphism of groups;
the affine-space factor has trivial fundamental group. Over $\C$
these statements also follow by deformation retraction and the
finite-cover comparison, and descend to characteristic zero.
The affine line has no nontrivial connected finite \'etale cover:
Riemann--Hurwitz applied to a hypothetical cover of $\mathbb P^1$
branched only at infinity gives
$2g-2=-2d+\sum(e_i-1)\leq-d-1$, impossible for $d>1$.
Its local group at infinity is nevertheless $\widehat{\Z}(1)$.
For the last example use coordinates $u=x/y$ and $v=z/y$ on the
complement; it is $\G_m\times\mathbb A^1$. In the local chart $z=1$,
the global coordinate is $u=x/y$, so its Kummer character has
valuations $(1,-1)$ on the two branches. The formula and its kernel
follow. The finite reductions of these maps, and the explicit failed
extension of a local cyclic cover, are checked in
\texttt{code/root\_checks.py}.
\end{proof}
